\documentclass[11pt]{article}

\usepackage{amsmath,amssymb,amsthm,mathtools}
\usepackage{geometry}
\usepackage{hyperref}
\usepackage{enumitem}
\usepackage{mathrsfs}

\geometry{margin=1in}

\title{Quantization as a Deformation of State Space:\\
Convex Geometry, Noncommutativity, and the Failure of Canonical Functoriality}

\author{}
\date{}

\newtheorem{theorem}{Theorem}
\newtheorem{definition}{Definition}
\newtheorem{proposition}{Proposition}
\newtheorem{remark}{Remark}
\newtheorem{corollary}{Corollary}

\begin{document}

\maketitle

\begin{abstract}
Canonical quantization is traditionally formulated as a homomorphism
from classical observables equipped with a Poisson bracket
to noncommuting operators satisfying canonical commutation relations.
It is well known that such a homomorphism cannot exist globally.
We argue that this obstruction is structural rather than technical.
Quantization is not a deformation of observables alone,
but a deformation of the convex geometry of states.
Classical mechanics is commutative probability on a simplex,
whereas quantum mechanics is noncommutative probability on a non-simplex convex set.
Noncommutativity, contextuality, and the failure of canonical covariance
follow directly from this geometric transition.
\end{abstract}

\tableofcontents

\section{Introduction}

Let $(M,\omega)$ be a classical phase space.
Canonical quantization is commonly expressed as a rule
\begin{equation}
\{f,g\} \longmapsto \frac{1}{i\hbar}[Q(f),Q(g)],
\end{equation}
where $\{\cdot,\cdot\}$ denotes the Poisson bracket.

It is known that no quantization map
\[
Q: C^\infty(M) \to \mathcal{B}(\mathcal{H})
\]
can preserve the Poisson structure, multiplicative structure,
and canonical covariance simultaneously for all observables.

Rather than treating this as an algebraic anomaly,
we propose a structural reformulation:

\begin{center}
\textbf{Quantization is a change of state space geometry.}
\end{center}

Observable algebras are secondary; state spaces are primary.

\section{Classical Mechanics as Commutative Probability}

\subsection{Algebraic structure}

Let
\[
\mathcal{A}_{\mathrm{cl}} = C_0(M)
\]
with pointwise multiplication.

States are positive normalized linear functionals:
\[
\mu(f) = \int_M f\, d\mu.
\]

\subsection{Convex structure}

The classical state space is
\[
\mathcal{S}_{\mathrm{cl}} = \mathcal{P}(M).
\]

\begin{proposition}
$\mathcal{S}_{\mathrm{cl}}$ is a Choquet simplex.
\end{proposition}

\begin{proof}
Extreme points are Dirac measures $\delta_x$.
Every probability measure admits a unique barycentric decomposition
in terms of these extreme points.
\end{proof}

\begin{remark}
The uniqueness of decomposition characterizes classical probability.
\end{remark}

Thus classical mechanics is equivalently:

\[
\text{Commutative } C^*\text{-algebra}
\quad \Longleftrightarrow \quad
\text{Simplex state space}.
\]

\section{Quantum Mechanics as Noncommutative Probability}

\subsection{Algebraic structure}

Let $\mathcal{H}$ be a Hilbert space.
Observables are elements of a noncommutative $C^*$-algebra
$\mathcal{A}_{\mathrm{qm}} \subset \mathcal{B}(\mathcal{H})$.

States are density operators:
\[
\mathcal{S}_{\mathrm{qm}} =
\{ \rho \ge 0 \mid \mathrm{Tr}(\rho)=1 \}.
\]

\subsection{Convex geometry}

\begin{proposition}
$\mathcal{S}_{\mathrm{qm}}$ is convex but not a simplex.
\end{proposition}

\begin{proof}
Mixed states admit infinitely many decompositions into pure states.
Uniqueness fails; hence not a simplex.
\end{proof}

\begin{remark}
Non-uniqueness of decomposition is equivalent to quantum interference.
\end{remark}

\section{Duality Between Observables and States}

For any $C^*$-algebra $\mathcal{A}$,
states are positive normalized linear functionals on $\mathcal{A}$.

Conversely, affine functionals on the state space reconstruct observables.

Thus:

\begin{center}
Changing the state space necessarily changes the observable algebra.
\end{center}

One cannot preserve classical state geometry while deforming
only the multiplication law.

\section{Quantization as Convex Deformation}

\begin{definition}
Quantization is the replacement
\[
\mathcal{S}_{\mathrm{cl}} = \mathcal{P}(M)
\quad \longrightarrow \quad
\mathcal{S}_{\mathrm{qm}} = \mathcal{D}(\mathcal{H}),
\]
where a simplex state space is replaced by a non-simplex convex set.
\end{definition}

Observables are then reconstructed as affine functionals
on the new state space.

\section{Geometric Obstruction to Canonical Quantization}

Assume a quantization map $Q$ preserving Lie structure.
Dualizing, one obtains an affine embedding of
$\mathcal{S}_{\mathrm{qm}}$ into $\mathcal{S}_{\mathrm{cl}}$.

\begin{theorem}
There is no affine isomorphism between a simplex and a non-simplex convex set.
\end{theorem}

\begin{proof}
The simplex property is invariant under affine isomorphism.
Quantum state space violates uniqueness of decomposition.
\end{proof}

\begin{corollary}
A global Poisson-to-commutator homomorphism cannot exist.
\end{corollary}

The obstruction is therefore geometric.

\section{Symmetries and Canonical Transformations}

Classically, a canonical transformation
\[
\phi: M \to M
\]
acts by pushforward on states:
\[
\mu \mapsto \phi_*\mu.
\]

Quantum mechanically, symmetries are automorphisms of the operator algebra,
implemented (when possible) by unitary operators.

Since the state space has changed,
there is no reason for every classical canonical transformation
to lift to a quantum symmetry.

Functoriality fails because the underlying convex geometry differs.

\section{Contextuality and Loss of Global Sample Space}

In classical theory:

\begin{itemize}
\item All observables are functions on one space $M$.
\item A single global event space exists.
\end{itemize}

In quantum theory:

\begin{itemize}
\item Only commuting observables admit joint spectra.
\item No single global sample space exists.
\end{itemize}

Quantization therefore represents the loss of a global Kolmogorovian event space.

\section{Relation to Deformation Quantization}

Deformation quantization replaces pointwise multiplication by
a star-product:
\[
f \star g = fg + \frac{i\hbar}{2}\{f,g\} + O(\hbar^2).
\]

However, even this formal deformation implicitly changes
the dual state space structure.

The obstruction to global canonical quantization
manifests as nontrivial higher-order terms in $\hbar$.

\section{Semiclassical Limit}

As $\hbar \to 0$:

\begin{itemize}
\item Commutators approximate Poisson brackets.
\item Quantum states concentrate toward measures.
\end{itemize}

The simplex structure is recovered only in the limit.
Thus classical mechanics appears as a singular convex-geometric limit.

\section{Conclusion}

Quantization is not a homomorphism of algebras.
It is a transformation of probability geometry:

\[
\text{Simplex state space}
\quad \longrightarrow \quad
\text{Non-simplex state space}.
\]

Noncommutativity,
contextuality,
and the failure of canonical covariance
follow naturally from this geometric shift.

The traditional quantization problem is unsolvable
because it attempts to preserve classical convex structure
while deforming algebraic relations.

The correct perspective is therefore:

\begin{center}
\textbf{Quantization is a deformation of state space,
not a map of observables.}
\end{center}

\end{document}
